\chapter{Trauma}

\section*{Patient focus}
After an accident or injury, imaging helps doctors quickly find broken bones, internal bleeding, and organ damage. A CT scan is often the most important test after major trauma. It is fast, shows almost every part of the body, and helps surgeons decide what needs urgent treatment. You may also have X-rays, ultrasound, or MRI depending on your injuries.


\label{ch:ch58_trauma}

\section{Initial Strategy}
Imaging follows resuscitation and the clinical question. The primary survey follows ATLS principles: life-threatening injuries are identified and treated before detailed imaging.

\subsection{Primary Survey Imaging}
\begin{itemize}
  \item \textbf{Chest X-ray}: pneumothorax, haemothorax, rib fractures, widened mediastinum
  \item \textbf{Pelvic X-ray}: open-book fracture, vertically unstable pelvis
  \item \textbf{FAST ultrasound}: free fluid in abdomen (perisplenic, perihepatic, pelvis, subxiphoid)
  \item \textbf{E-FAST}: adds bilateral anterior chest for pneumothorax
\end{itemize}

\subsection{Definitive Imaging}
\begin{itemize}
  \item \textbf{CT with IV contrast}: cornerstone of polytrauma imaging in haemodynamically stable patients
  \item \textbf{CT angiography}: suspected vascular injury, active hemorrhage
  \item \textbf{MRI}: spinal cord injury, ligamentous instability, occult fractures, diffuse axonal injury
  \item \textbf{CT head}: always non-contrast for acute trauma
\end{itemize}

\section{Trauma CT Protocol}
A standard polytrauma CT protocol includes:

\begin{table}[h]
\centering
\caption{Standard polytrauma CT protocol}
\begin{tabular}{|l|l|l|}
\hline
\textbf{Body region} & \textbf{Technique} & \textbf{Key findings} \\
\hline
Abdomen & CECT arterial + portal venous & Organ injury, active bleed \\
C Spine & NCCT & Fracture, malalignment, subluxation \\
Chest & CECT & Haemothorax, pneumothorax, aortic injury \\
Head & NCCT & Hemorrhage, midline shift, fractures \\
Pelvis & CECT & Fracture, bladder injury, vascular injury \\
\hline
\end{tabular}
\end{table}

\textbf{Timing}: Arterial phase (25--30 s) for active hemorrhage and vascular injury; portal venous phase (60--70 s) for organ parenchymal injury. Delayed phase (5--10 min) only if urinary tract injury is suspected.

\section{Interpretation---Systematic Approach}
Use the \textbf{ABCD} trauma search pattern on every scan:

\begin{enumerate}
  \item \textbf{A---Airway}: pneumomediastinum, tracheal deviation, foreign body
  \item \textbf{B---Breathing}: pneumothorax (anterior on supine CT), haemothorax, pulmonary contusion
  \item \textbf{C---Circulation}: active contrast extravasation (blush), pseudoaneurysm, AV fistula, free fluid
  \item \textbf{D---Disability}: intracranial hemorrhage, axonal injury, spinal cord compression
\end{enumerate}

Then assess: fractures (bone windows), visceral organ injury, hollow viscus injury (free air, wall thickening), retroperitoneal hemorrhage, and iatrogenic devices.

\section{AAST Organ Injury Grading}
The American Association for the Surgery of Trauma (AAST) grading systems are used to classify organ injuries and guide management:

\begin{table}[h]
\centering
\caption{AAST Spleen Injury Scale (selected grades)}
\begin{tabular}{|p{0.18\textwidth}|p{0.70\textwidth}|}
\hline
\textbf{Grade} & \textbf{Description} \\
\hline
I & Haematoma $<$10\% surface area; laceration $<$1 cm depth \\
II & Haematoma 10--50\%; laceration 1--3 cm, not in hilar vessels \\
III & Haematoma $>$50\% or expanding; laceration $>$3 cm or involving trabecular vessels \\
IV & Laceration involving segmental or hilar vessels with $>$25\% devascularization \\
V & Shattered splein or hilar vascular injury with complete devascularization \\
\hline
\end{tabular}
\end{table}

\begin{pearlbox}[Grading and management]
AAST grade alone does not determine management. Haemodynamic status, patient age, comorbidities, associated injuries, and institutional capability all influence the decision between observation, angiographic embolization, and surgery. Higher-grade injuries in haemodynamically unstable patients still require immediate intervention.
\end{pearlbox}

\section{When to Recommend Angiography vs Surgery}

\begin{table}[h]
\centering
\caption{Decision pathway for hemorrhage control}
\begin{tabular}{|p{0.35\textwidth}|p{0.28\textwidth}|p{0.23\textwidth}|}
\hline
\textbf{Scenario} & \textbf{Angiography} & \textbf{Surgery} \\
\hline
Active arterial blush, stable & Embolization & --- \\
Active arterial blush, unstable & --- & Emergency laparotomy \\
Shattered organ, vascular pedicle & --- & Surgery \\
Solid organ injury, no active bleed & Observation/angio & --- \\
Solid organ injury, venous bleeding & --- & Surgery \\
Suspected pelvic fracture bleed & Angiography + embolization & External fixation \\
\hline
\end{tabular}
\end{table}

\section{Delayed Complications}
Recognize these delayed complications on follow-up imaging:

\begin{itemize}
  \item \textbf{Pseudoaneurysm}: focal outpouching at arterial injury site, shows filling on arterial phase---treat with embolization
  \item \textbf{Arteriovenous fistula}: early venous filling during arterial phase; may cause high-output cardiac failure
  \item \textbf{Post-traumatic infarction}: wedge-shaped non-enhancing area in solid organ
  \item \textbf{Biloma/biliary leak}: fluid collection near liver, may need ERCP or percutaneous drainage
  \item \textbf{Pancreatic duct injury}: initially subtle; delayed CT may show peripancreatic fluid, pseudocyst
  \item \textbf{Bowel perforation}: may be delayed 24--48 h; serial imaging with water-soluble contrast if initial study equivocal
  \item \textbf{Spinal cord injury}: delayed myelomalacia, syrinx, cord atrophy on follow-up MRI
\end{itemize}

\section{Trauma in Special Populations}
\begin{itemize}
  \item \textbf{Pediatric}: solid organ injury more common than skeletal; FAST is useful but may miss retroperitoneal injury; radiation dose reduction is essential
  \item \textbf{Elderly}: falls from standing height can cause serious injury; occult hip fractures need MRI if X-ray normal; brain atrophy masks intracranial haematoma
  \item \textbf{Anticoagulated patients}: higher risk of delayed intracranial hemorrhage; repeat CT head at 6--24 h if initial scan is normal but clinical concern persists
\end{itemize}

\section{Communication}
Do not delay life-saving treatment for imaging. Communicate the following immediately:
\begin{itemize}
  \item Active hemorrhage or contrast blush
  \item Tension physiology (tension pneumothorax, cardiac tamponade)
  \item Vascular injury requiring intervention
  \item Unstable spinal fracture with cord compression
  \item Bowel perforation or mesenteric injury
  \item Compartment syndrome
\end{itemize}

\begin{pearlbox}[Trauma communication]
Use structured communication (SBAR). State the finding, its clinical significance, and whether it requires immediate intervention. ``Active splenic hemorrhage with blush, Grade III splenic laceration, haemodynamically stable---consider angiography'' is more useful than a purely descriptive report.
\end{pearlbox}

\clearpage
